Since $A$ is Noetherian, $\mathfrak p_i^{r_i}\subseteq\mathfrak q_i$ for some $r_i$. Localizing at $\mathfrak p_i$ and contracting gives $\mathfrak p_i^{(r)}\subseteq\mathfrak q_i$ for $r\geq r_i$, since a $\mathfrak p_i$-primary ideal is contracted from $A_{\mathfrak p_i}$.

If $\mathfrak p_i$ is isolated, it is minimal. The Noetherian local ring $A_{\mathfrak p_i}$ has dimension zero and is Artinian, so its maximal ideal has a zero power. Every other primary component becomes the unit ideal upon localization. Hence $\mathfrak q_iA_{\mathfrak p_i}=0$ and, for large $r$, contraction gives $\mathfrak q_i=\mathfrak p_i^{(r)}$.

If $\mathfrak p_i$ is embedded, $A_{\mathfrak p_i}$ has positive dimension and is not Artinian. Its maximal-ideal powers are all distinct: equality of two consecutive powers would, by Nakayama's lemma, make that power zero and the ring Artinian. Their contractions $\mathfrak p_i^{(r)}$ are therefore distinct. For $r\geq r_i$, replacing $\mathfrak q_i$ by $\mathfrak p_i^{(r)}$ still gives intersection zero. The decomposition remains minimal, since otherwise deleting a redundant component would contradict the first uniqueness theorem for the associated primes of $(0)$.
